\begin{abstract}
Large language model (LLM) agents that interact with external tools and data sources are vulnerable to \emph{indirect prompt injection} (IPI), where adversarial instructions embedded in untrusted content hijack agent behavior.
Despite growing interest in IPI defenses, there is no comprehensive benchmark that evaluates diverse defense strategies under controlled, reproducible conditions.
We introduce the \textbf{Agent Security Sandbox} (\asb{}), a modular benchmark framework for systematically evaluating IPI defenses.
\asb{} comprises 565 test cases (352 attack, 213 benign) spanning 6 attack types and 54 injection techniques, a unified evaluation pipeline with rule-based judging, and implementations of 11 defense strategies (D0--D10) covering prompt-level, tool-gating, content-level, and multi-signal approaches.
We conduct a large-scale evaluation across 4 frontier LLMs (GPT-4o, Claude~4.5 Sonnet, DeepSeek~V3.1, Gemini~2.5 Flash) with 132 evaluation runs, including ablation studies, defense composition analysis, and adaptive attack testing.
Our evaluation reveals four key findings: (1)~simple prompt-level defenses (sandwich prompting) reduce attack success by 97.5\% while preserving 93.4\% benign task completion, outperforming complex tool-gating approaches; (2)~model robustness varies by 4--5$\times$ across families, with Claude exhibiting strong inherent resistance; (3)~defense composition yields additive but not super-additive security gains; (4)~input sanitization is \emph{counterproductive}---it removes adversarial cues that help aligned models self-detect injections, increasing attack success by 4\%.
We also contribute \textbf{Contextual Integrity Verification} (\civ{}), a novel multi-signal defense combining prompt-layer framing with risk-stratified tool gating informed by contextual integrity theory.
An ablation-driven iterative design (CIV~v1$\to$v2) replaces a harmful counterfactual heuristic with plan deviation detection and upgrades keyword-based tool affinity to embedding-based semantic compatibility, achieving 78\% attack reduction with a mechanism orthogonal to prompt-level approaches.
We release \asb{} as an open-source toolkit for reproducible agent security research.\footnote{Code and data: \url{https://github.com/X-PG13/agent-security-sandbox}}
\end{abstract}
